% SEGMENT OLP-0091-S01
% Part: first-order-logic
% Chapter: natural-deduction
% Section: proof-theoretic-notions

\documentclass[../../../include/open-logic-section]{subfiles}

\begin{document}

\iftag{FOL}
      {\olfileid{fol}{ntd}{ptn}}
      {\olfileid{pl}{ntd}{ptn}}

\olsection{நிறுவல்-கோட்பாட்டுக் கருத்துகள்}

\begin{editorial}
இந்தப் பிரிவு இயல்பான வருவித்தலுக்கான வருவிக்கத்தக்கதன்மை உறவு மற்றும்
முரணின்மை ஆகியவற்றின் வரையறைகளைத் தொகுக்கிறது.
\end{editorial}

\begin{explain}
பல முக்கியமான பொருண்மைக் கருத்துகளை (செல்லுபடித்தன்மை, பின்விளைவு,
நிறைவுறத்தக்கதன்மை) நாம் வரையறுத்ததைப் போலவே, அவற்றுடன் தொடர்புடைய
\emph{நிறுவல்-கோட்பாட்டுக் கருத்துகளை} இப்போது வரையறுக்கிறோம். இவை
!!{structure}s[loc] !!{sentence}s நிறைவுறுவதைச் சார்ந்து வரையறுக்கப்படாமல்,
குறிப்பிட்ட !!{sentence}s மற்றவற்றிலிருந்து !!{derivability} அல்லது
!!{nonderivability} உடையனவாக இருப்பதைச் சார்ந்து வரையறுக்கப்படுகின்றன.
இந்தக் கருத்துகள் ஒன்றுபடுகின்றன என்பது ஒரு முக்கியமான கண்டுபிடிப்பு. அவை
ஒன்றுபடுவதே \emph{சரித்தன்மைத் தேற்றம்} மற்றும் \emph{நிறைவுத் தேற்றம்}
ஆகியவற்றின் உள்ளடக்கம்.
\end{explain}

% SEGMENT OLP-0091-S02
\begin{defn}[தேற்றங்கள்]
ஒவ்வோர் எடுகோளும் !!{discharged} ஆகும் வகையில், இயல்பான வருவித்தலில்
!!{sentence}~$!A$ இன் !!a{derivation} இருந்தால், $!A$ ஒரு \emph{தேற்றம்}.
$!A$ தேற்றமாக இருந்தால் $
\Proves !A$ என்றும், இல்லையெனில் $\Proves/ !A$ என்றும் எழுதுகிறோம்.
\end{defn}

% SEGMENT OLP-0091-S03
\begin{defn}[!!^{derivability}]
!!{sentence}s கொண்ட~$\Gamma$ கணத்திலிருந்து !!^a{sentence}~$!A$ ஆனது
\emph{!!{derivable}} என்பதை $\Gamma \Proves !A$ என்று எழுதுகிறோம். இதன்
பொருள், முடிவு~$!A$ ஐக் கொண்டதும், ஒவ்வோர் எடுகோளும் !!{discharged} ஆகவோ
$\Gamma$ வில் இருக்கவோ செய்கின்றதுமான !!a{derivation} ஒன்று உள்ளது என்பதாகும்.
$!A$ ஆனது~$\Gamma$ இலிருந்து !!{derivable} அல்ல எனில் $\Gamma \Proves/ !A$
என்று எழுதுகிறோம்.
\end{defn}

% SEGMENT OLP-0091-S04
\begin{defn}[முரணின்மை]
!!{sentence}s கொண்ட~$\Gamma$ கணம், $\Gamma \Proves \lfalse$ எனில், அப்போதும்
அப்போது மட்டுமே \emph{முரணுடையது}. $\Gamma$ முரணுடையதல்ல எனில், அதாவது
$\Gamma \Proves/ \lfalse$ எனில், அது \emph{முரணற்றது} என்கிறோம்.
\end{defn}

% SEGMENT OLP-0091-S05
\begin{prop}[தற்சுட்டுப் பண்பு]
\ollabel{prop:reflexivity}
$!A \in \Gamma$ எனில், $\Gamma \Proves !A$.
\end{prop}

\begin{proof}
எடுகோள்~$!A$ தனியே~$!A$ இன் !!a{derivation}; அதன் ஒவ்வொரு
!!{undischarged} எடுகோளும் (அதாவது $!A$)~$\Gamma$ வில் உள்ளது.
\end{proof}
  
% SEGMENT OLP-0091-S06
\begin{prop}[ஒருபோக்குத்தன்மை]
\ollabel{prop:monotonicity}
$\Gamma \subseteq \Delta$ மற்றும் $\Gamma \Proves !A$ எனில், $\Delta
\Proves !A$.
\end{prop}

\begin{proof}
$\Gamma$ இலிருந்து $!A$ இன் எந்த !!{derivation} உம்~$\Delta$ இலிருந்து
$!A$ இன் !!a{derivation} உம் ஆகும்.
\end{proof}

% SEGMENT OLP-0091-S07
\begin{prop}[கடப்புப் பண்பு]
\ollabel{prop:transitivity}
$\Gamma \Proves !A$ மற்றும் $\{!A\} \cup \Delta \Proves
!B$ எனில், $\Gamma \cup \Delta \Proves !B$.
\end{prop}

\begin{proof}
$\Gamma \Proves !A$ எனில், ஒவ்வொரு !!{undischarged} எடுகோளும்~$\Gamma$ வில்
இருக்குமாறு $!A$ இன் !!a{derivation}~$\delta_0$ உள்ளது. $\{!A\} \cup
\Delta \Proves !B$ எனில், ஒவ்வொரு !!{undischarged} எடுகோளும்
$\{!A\} \cup \Delta$ வில் இருக்குமாறு $!B$ இன்
!!a{derivation}~$\delta_1$ உள்ளது. இப்போது பின்வருவதைக் கருதுக:
\begin{prooftree}
  \AxiomC{$\Delta, \Discharge{!A}{1}$}
  \RightLabel{$\delta_1$}
  \DeduceC{$!B$}
  \DischargeRule{\Intro{\lif}}{1}
  \UnaryInfC{$!A \lif !B$}
  \AxiomC{$\Gamma$}
  \RightLabel{$\delta_0$}
  \DeduceC{$!A$}
  \RightLabel{\Elim{\lif}}
  \BinaryInfC{$!B$}
\end{prooftree}
இப்போது !!{undischarged} எடுகோள்கள் அனைத்தும் $\Gamma \cup \Delta$ வில்
உள்ளன; ஆகவே இது $\Gamma \cup \Delta \Proves !B$ என்பதைக் காட்டுகிறது.
\end{proof}

% SEGMENT OLP-0091-S08
$\Gamma = \{!A_1, !A_2, \ldots, !A_k\}$ ஒரு முடிவுறு கணமாக இருக்கும்போது,
$\Gamma \Proves !B$ என்பதற்கு எளிய குறியீடாக $!A_1,!A_2,\ldots,!A_k
\Proves !B$ ஐப் பயன்படுத்தலாம். குறிப்பாக $!A \Proves !B$ என்பது
$\{!A\} \Proves !B$ என்பதைக் குறிக்கும்.

$\Gamma \Proves !A$ மற்றும் $!A \Proves !B$ எனில், $\Gamma \Proves !B$
என்பதைக் கவனிக்கவும். மேலும் $!A_1, \dots, !A_n \Proves !B$ என்றும்,
ஒவ்வொரு~$i$ க்கும் $\Gamma \Proves !A_i$ என்றும் இருந்தால், $\Gamma
\Proves !B$ என்பதும் பின்வருகிறது.

% SEGMENT OLP-0091-S09
\begin{prop}
\ollabel{prop:incons}
பின்வருவன சமானமானவை.
    \begin{enumerate}
      \item \( \Gamma \) முரணுடையது.
      \item ஒவ்வொரு !!{sentence}~\( {!A} \) க்கும் \( \Gamma \Proves {!A} \).
      \item ஏதோவொரு !!{sentence}~\( {!A} \) க்கு \( \Gamma \Proves {!A} \)
        மற்றும் \( \Gamma \Proves \lnot {!A} \).
    \end{enumerate}
\end{prop}

\begin{proof}
பயிற்சி.
\end{proof}

\tagprob{FOL}
\begin{prob}
\olref[fol][ntd][ptn]{prop:incons} ஐ நிறுவுக.
\end{prob}
\tagendprob

\tagprob{notFOL}
\begin{prob}
\olref[pl][ntd][ptn]{prop:incons} ஐ நிறுவுக.
\end{prob}
\tagendprob

% SEGMENT OLP-0091-S10
\begin{prop}[இறுக்கத்தன்மை]
  \ollabel{prop:proves-compact}
  \begin{enumerate}
  \item $\Gamma \Proves !A$ எனில், $\Gamma_0 \subseteq \Gamma$ என்ற ஒரு
    முடிவுறு உட்கணம் இருந்து, $\Gamma_0 \Proves !A$.
  \item $\Gamma$ வின் ஒவ்வொரு முடிவுறு உட்கணமும் முரணற்றது எனில்,
    $\Gamma$ முரணற்றது.
  \end{enumerate}
\end{prop}

\begin{proof}
  \begin{enumerate}
    \item $\Gamma \Proves !A$ எனில், $\Gamma$ இலிருந்து~$!A$ இன்
      !!a{derivation}~$\delta$ உள்ளது. $\delta$ வின் !!{undischarged}
      எடுகோள்களின் கணத்தை~$\Gamma_0$ என்க. எந்த !!{derivation} உம்
      முடிவுறுவதால், $\Gamma_0$ வில் முடிவுறு எண்ணிக்கையிலான
      !!{sentence}s மட்டுமே இருக்க முடியும். ஆகவே $\delta$ ஆனது
      $\Gamma_0 \subseteq \Gamma$ என்ற முடிவுறு கணத்திலிருந்து~$!A$ இன்
      !!a{derivation}.
    \item இது $!A \ident \lfalse$ என்ற சிறப்புநிலையில் (1) இன்
      எதிர்மறைநேர்மாற்று.
  \end{enumerate}
\end{proof}

\end{document}
